\section{Przebieg gry}
\subsection{Ogólna zasada}
Gra \textit{Labirynt śmierci} polega na wykonywaniu przez graczy określonych czynności w ściśle określonej kolejności.
Jej celem jest odnalezienie przez oddział i zniszczenie magicznych Czarnych Wrót.

Decyzje dotyczące zachowania się postaci podczas gry mogą być podejmowane wspólnie przez wszystkich graczy lub też, jeżeli tak zostanie zdecydowane, każdy z graczy może podejmować decyzje indywidualnie.

\textit{Labirynt śmierci} rozgrywany jest w etapach.
Każdy etap jest to blok czynności, które gracze powinni wykonać przed przystąpieniem do następnego etapu.
Czynności te muszą być wykonywane w przypisanem im niżej kolejności:

\subsection{Umieszczenie nowego segmentu}
\begin{enumerate}
    \item
        Oddział ustala drogę, którą opuści zajmowany obecnie segment, czyli kierunek, w którym pójdzie.
    \item
        Jeden z graczy losowy wybiera z pojemnika nowy segment labiryntu (żeton).
    \item
        Nowy segment układany jest obok tego, który oddział opuszcza, na przedłużeniu drogi, którą zdecydował się on wybrać.
        Nowy segment musi zostać ułożony tak, by pasował do segmentu opuszczanego, to znaczy: korytarz musi stanowić przedłużenie korytarza, a przejście musi łączyć się z przejściem.
        Od decyzji graczy zależy w jaki sposób zostanie ułożony taki segment, w przypadku którego możliwe jest więcej niż jedno prawidłowe ułożenie.
\end{enumerate}

\subsection{Poszukiwanie pułapek}
\begin{enumerate}
    \item
        Jeden z graczy rzuca 1K6, by sprawdzic czy drzwi, przez które oddział ma zamiar przejść zawierają pułapkę.
        Poszukiwanie pułapek (rzucanie kostką) odbywa się tylkoe wtedy, gdy oddział chce wejść do nieodwiedzanej dotychczas komnaty, bądź też opuszcza komnątę inną drogą, niż do niej wszedł.
        A więc oddział nie szuka pułapek wtedy, gdy porusza się korytarzem.
        Drzwi zawierają pułapkę, jeśli wyrzucone zostanie 1 oczko.
    \item
        Po odnalezieniu połapki (wyrzucone 1 oczko) należy wydelegować jedną postać do jej zbadania.
        Delegowana postać powinna posiadać w grupie zdolności hasło \textit{pułapki}, gdyż tylko wtedy może podjąć próbe unieszkodliwienia znalezionej pułapki (patrz 7.0 NUMER).
    \item
        Jeżeli próba unieszkodliwienia pułapki zakończy się niepowodzeniem lub w ogóle nie zostanie podjęca, uważa się, że pułapka zadziałała.
        Gracze muszą ustalić jej rodzaj rzucając 1K6 i odczytując rezultat z tabeli 7.1 Pułapki NUMER!
\end{enumerate}

\subsection{Spotkanie z potworem}
\begin{enumerate}
    \item
        Oddział wykonuje ruch.
        Odbywa się to przez umieszczenie jego żetornu na nowym segmencie.
        Następną czynnością jest sprawdzenie, czy w nowym segmencie znajdują się potwory.
        Jeden z graczy rzuca 1K6.
        Spotkanie z potworem następuje przy wyrzuceniu:
        \begin{enumerate}
            \item
                1, 2 lub 3 oczek, jeżeli nowy segment jest niezbadaną dotychczas komnatą.
            \item
                1 oczka jeżeli nowy segment jest fragmentem korytarza lub komnatą, którą oddział juz poprzednio odwiedził
        \end{enumerate}
    \item
        Oddział nie musi od razu przystępować do walki z napotkanym potworem.
        Może przedtem podjąć próbę negocjacji lub wykupienia się (patrz 8.0 NUMER)
    \item
        Jeżeli próby negocjacju lub wykupienia się zawiodą, lub nie zostaną podjęte, rozpoczyna się walka.
        Składa się ona z następujących faz (czynności):
        \begin{enumerate}
            \item
                atak oddziału,
            \item
                atak potwora,
            \item
                atak Czarnych Wrót,
            \item
                reorganizacja szyku oddziału,
            \item
                reorganizacja szyku potworów.
        \end{enumerate}
        Powyższe fazy muszą być wykonywane w podanej kolejności.
    \item
        Po zabiciu potwora (lub wszystkich potworów stanowiących napotkaną grupę) oddział ma prawo odebrać mu posiadane przez niego skarby.
        Ustala się ile i jakie skarby potwór posiadał (patrz 14.0 NUMER), a następnie dzieli się je między członków oddziału.
    \item
        Jeżeli potwór zostanie zabity, pozostali przy życiu członkowie oddziału zdobywają określoną liczbę punktów doświadczenia (patrz 12.0 NUMER).
\end{enumerate}

\subsection{Badanie znalezisk}
Znaleziska są to rzeczy znalezione przez oddział w komnatach (tylko).
\begin{enumerate}
    \item
        Jedna z postaci jest delegowana do zbadania znaleziska.
    \item
        Badająca postać (gracz) rzuca 1K6.
        Z tabeli 13.9 \textit{Znaleziska} NUMER, z kolumny odpowiadającej jego rodzajowi odczytywana jest dokładna nazwa tego, co zostało znalezione.
    \item
        Wprowadzane są w życie wszelkie efekty, wiążące się z danym znaleziskiem.
\end{enumerate}
Wszelkie czynności wymienione w punktach A do D NUMER! nazywane są wspólnie etapem gry.
Po ich wykonaniu gracze przystępują do kolejnego etapu, czyli od początku do punktu A NUMER!.
